\documentclass[../../../include/open-logic-section]{subfiles}

\begin{document}
\olfileid[id]{sth}{ord-arithmetic}{intro}

\olsection{Pendahuluan}

Dalam \olref[ordinals][]{chap}, kita mengembangkan suatu teori bilangan
ordinal. Dalam \olref[spine][]{chap}, kita melihat bahwa ordinal dapat
dipandang sebagai tulang punggung yang menjadi kerangka konstruksi bagian
hierarki selebihnya. Namun, itu bukan satu-satunya peran ordinal. Ada pula
tugas melakukan aritmetika ordinal.

Kita telah mengisyaratkan hal ini sebelumnya dalam
\olref[ordinals][idea]{sec}, ketika membicarakan $\omega$, $\omega+1$, dan
$\omega+\omega$. Saat itu kita berbicara secara informal; sekarang tiba
waktunya untuk menjabarkannya dengan tepat. Namun, perlu disebutkan bahwa bab
ini tidak banyak memuat filsafat: hanya perkembangan teknis, disertai
pengamatan yang (agak) menarik bahwa kita dapat melakukan hal yang sama dengan
dua cara berbeda.

\end{document}
